\ifdefined\gkver\else\def\gkver{student}\fi
\documentclass[\gkver]{gaokaozhenti}
\begin{document}

\chapter{2004年全国理综4}
%% paperType: 理综物理部分

\begin{enumerate}
\item
%% number: 14
%% paperName: 2004年全国理综4第 14 题 6 分
%% typeId: 1
%% type: 单选题
%% chapter: 原子和原子核
%% point: 核裂变
%% method: 无
%% score: 6
%% degree: 850
%% duplicateId: 0
%% body:
在核反应方程式 $\ce{^{235}_{92}U} + \ce{^{1}_{0}n} \to \ce{^{90}_{38}Sr} + \ce{^{136}_{54}Xe} + k\ce{X}$ 中\xzanswer{B}
\fourchoices[answer=B]
{X 是中子，$k=9$}
{X 是中子，$k=10$}
{X 是质子，$k=9$}
{X 是质子，$k=10$}
%% answer: B
%% memo:
\memoanswer{
核反应方程满足质量数和电荷数守恒。由裂变规律可知，X 是中子 $\ce{^{1}_{0}n}$，由质量数守恒得 $235+1=90+136+k$，可得 $k=10$，故 B 正确。
}

\item
%% number: 15
%% paperName: 2004年全国理综4第 15 题 6 分
%% typeId: 1
%% type: 单选题
%% chapter: 电磁感应
%% point: 电磁感应中图象
%% method: 无
%% score: 6
%% degree: 650
%% duplicateId: 0
%% body:
如图所示，在 $x\leq 0$ 的区域内存在匀强磁场，磁场的方向垂直于 $xy$ 平面（纸面）向里。具有一定电阻的矩形线框 abcd 位于 $xy$ 平面内，线框的 ab 边与 $y$ 轴重合。令线框从 $t=0$ 的时刻起由静止开始沿 $x$ 轴正方向做匀加速运动，则线框中的感应电流 $I$（取逆时针方向的电流为正）随时间 $t$ 的变化图线 $I-t$ 图可能是下图中的哪一个？\xzanswer{D}

\onepicture[w=5.0cm]{figs/15a.png}

\fourchoices[answer=D, ispicture=true, h=1.9cm]
{figs/15b.png}
{figs/15c.png}
{figs/15d.png}
{figs/15e.png}
%% answer: D
%% memo:
\memoanswer{
cd 边切割磁感线产生的电动势为 $E=Blv$，又 $I=\frac{E}{R}=\frac{Blv}{R}=\frac{Blat}{R}$，与 $t$ 成正比，由楞次定律（或右手定则）可知感应电流为顺时针，故 D 正确。
}

\item
%% number: 16
%% paperName: 2004年全国理综4第 16 题 6 分
%% typeId: 1
%% type: 单选题
%% chapter: 气体性质
%% point: 热力学第一定律
%% method: 无
%% score: 6
%% degree: 650
%% duplicateId: 0
%% body:
一定质量的理想气体，从某一状态开始，经过系列变化后又回一开始的状态，用 $W_{1}$ 表示外界对气体做的功，$W_{2}$ 表示气体对外界做的功，$Q_{1}$ 表示气体吸收的热量，$Q_{2}$ 表示气体放出的热量，则在整个过程中一定有\xzanswer{A}
\fourchoices[answer=A]
{$Q_{1}-Q_{2}=W_{2}-W_{1}$}
{$Q_{1}=Q_{2}$}
{$W_{1}=W_{2}$}
{$Q_{1}>Q_{2}$}
%% answer: A
%% memo:
\memoanswer{
理想气体初态和末态相同，则温度相同，内能变化为零，由 $\Delta U=W+Q=W_{1}-W_{2}+Q_{1}-Q_{2}=0$ 可得 $Q_{1}-Q_{2}=W_{2}-W_{1}$，故 A 正确。
}

\item
%% number: 17
%% paperName: 2004年全国理综4第 17 题 6 分
%% typeId: 1
%% type: 单选题
%% chapter: 光学
%% point: 平面镜成像
%% method: 无
%% score: 6
%% degree: 650
%% duplicateId: 0
%% body:
图中 M 是竖直放置的平面镜，镜离地面的距离可调节。甲、乙二人站在镜前，乙离镜的距离为甲离镜的距离的 2 倍，如图所示。二人略错开，以便甲能看到乙的像。以 $l$ 表示镜的长度，$h$ 表示乙的身高，为使甲能看到镜中乙的全身像，$l$ 的最小值为\xzanswer{A}

\onepicture[h=3.4cm]{figs/17.png}

\fourchoices[answer=A]
{$\frac{1}{3}h$}
{$\frac{1}{2}h$}
{$\frac{3}{4}h$}
{$h$}
%% answer: A
%% memo:
\memoanswer{
如图所示，由对称性做出乙的像 CD，镜的最小长度即为图中 AB 的长度，由 $\Delta OAB\sim\Delta OCD$ 有 $\frac{AB}{CD}=\frac{1}{3}$，即 $AB=\frac{1}{3}h$，故 A 正确。

\onepicture[w=7.5cm]{figs/17_an.jpg}
}

\item
%% number: 18
%% paperName: 2004年全国理综4第 18 题 6 分
%% typeId: 1
%% type: 单选题
%% chapter: 机械波
%% point: 波的图象
%% method: 无
%% score: 6
%% degree: 750
%% duplicateId: 0
%% body:
已知：一简谐横波在某一时刻的波形图如图所示，图中位于 a、b 两处的质元经过四分之一周期后分别运动到 a$'$、b$'$ 处。某人据此做出如下判断：①可知波的周期，②可知波的传播速度，③可知的波的传播方向，④可知波的波长。其中正确的是\xzanswer{C}

\onepicture[h=3.0cm]{figs/18.png}

\fourchoices[answer=C]
{①和④}
{②和④}
{③和④}
{②和③}
%% answer: C
%% memo:
\memoanswer{
由图像可知波长 $\lambda=4\Um$，由 b 点振动方向可知波向右传播，但周期不知，也不能求得波速，故 C 正确。
}

\item
%% number: 19
%% paperName: 2004年全国理综4第 19 题 6 分
%% typeId: 1
%% type: 单选题
%% chapter: 运动和力
%% point: 牛顿第二定律
%% method: 整体与隔离
%% score: 6
%% degree: 550
%% duplicateId: 0
%% body:
如图，在倾角为 $\alpha$ 的固定光滑斜面上，有一用绳子拴着的长木板，木板上站着一只猫。已知木板的质量是猫的质量的 2 倍。当绳子突然断开时，猫立即沿着板向上跑，以保持其相对斜面的位置不变。则此时木板沿斜面下滑的加速度为\xzanswer{C}

\onepicture[h=2.3cm]{figs/19.png}

\fourchoices[answer=C]
{$\frac{1}{2}g\sin\alpha$}
{$g\sin\alpha$}
{$\frac{3}{2}g\sin\alpha$}
{$2g\sin\alpha$}
%% answer: C
%% memo:
\memoanswer{
对猫和木板整体受力分析，因猫位置不变，其加速度为零，由牛顿第二定律得 $3mg\sin\alpha=2ma$，解得 $a=\frac{3}{2}g\sin\alpha$，故 C 正确。
}

\item
%% number: 20
%% paperName: 2004年全国理综4第 20 题 6 分
%% typeId: 1
%% type: 单选题
%% chapter: 曲线运动
%% point: 竖直平面圆周运动
%% method: 无
%% score: 6
%% degree: 550
%% duplicateId: 0
%% body:
如图所示，轻杆的一端有一个小球，另一端有光滑的固定轴 O。现给球一初速度，使球和杆一起绕 O 轴在竖直面内转动，不计空气阻力，用 $F$ 表示球到达最高点时杆对小球的作用力，则 $F$\xzanswer{D}

\onepicture[h=3.2cm]{figs/20.png}

\fourchoices[answer=D]
{一定是拉力}
{一定是推力}
{一定等于 0}
{可能是拉力，可能是推力，也可能等于 0}
%% answer: D
%% memo:
\memoanswer{
球能通过最高点的条件是 $v>0$。①当最高点时杆对球作用力为零时，由 $mg=m\frac{v^{2}}{R}$ 得 $v=\sqrt{gR}$；②当 $0<v<\sqrt{gR}$ 时，杆对球是推力；③当 $v>\sqrt{gR}$ 时，杆对球是拉力，故 D 正确。
}

\item
%% number: 21
%% paperName: 2004年全国理综4第 21 题 6 分
%% typeId: 1
%% type: 单选题
%% chapter: 电场
%% point: 电势能电势差电场力做功
%% method: 无
%% score: 6
%% degree: 350
%% duplicateId: 0
%% body:
一平行板电容器的电容为 $C$，两板间的距离为 $d$，上板带正电，电量为 $Q$，下板带负电，电量也为 $Q$，它们产生的电场在很远处的电势为零。两个带异号电荷的小球用一绝缘刚性杆相连，小球的电量都为 $q$，杆长为 $l$，且 $l<d$。现将它们从很远处移到电容器内两板之间，处于图示的静止状态（杆与板面垂直），在此过程中电场力对两个小球所做总功的大小等于多少？（设两球移动过程中极板上电荷分布情况不变）\xzanswer{A}

\onepicture[h=2.6cm]{figs/21.png}

\fourchoices[answer=A]
{$\frac{Qlq}{Cd}$}
{$0$}
{$\frac{Qq(d-l)}{Cd}$}
{$\frac{Clq}{Qd}$}
%% answer: A
%% memo:
\memoanswer{
当两个带等量异号电荷的小球从很远处移到电场中同一点时，电场力做功为零，再将某一球沿电场线移到 $l$ 距离远，此时电场力做功即为所求。场强 $E=\frac{U}{d}=\frac{Q}{Cd}$，电场力做功 $W=qEl=\frac{Qql}{Cd}$，故 A 正确。
}

\item
%% number: 22
%% paperName: 2004年全国理综4第 22 题 4 分
%% typeId: 4
%% type: 实验题
%% chapter: 直线运动
%% point: 游标卡尺和螺旋测微仪
%% method: 无
%% score: 4
%% degree: 850
%% duplicateId: 0
%% body:
图中给出的是螺旋测微器测量一金属板厚度时的示数，读数应为\tkanswer[2.5cm]{7.324}$\Umm$（7.323$\sim$7.325 均给分）。
\onepicture[align=right, w=5.0cm]{figs/22a.jpg}
%% answer: 7.324（7.323～7.325 均给分）
%% memo:
\memoanswer{
固定读数为 $7\Umm$，可动读数为 $32.4\times 0.01\Umm=0.324\Umm$，故读数为 $7\Umm+0.324\Umm=7.324\Umm$。
}

\item
%% number: 22
%% paperName: 2004年全国理综4第 22 题 14 分
%% typeId: 4
%% type: 实验题
%% chapter: 电场
%% point: 描绘等势面
%% method: 无
%% score: 14
%% degree: 550
%% duplicateId: 0
%% body:
图中给出器材为：
电源 $E$（电动势为 $12\UV$，内阻不计），
木板 N（板上从下往上依次叠放白纸、复写纸、导电纸各一张），两个金属条 A、B（平行放置在导电纸上，与导电纸接触良好，用作电极），
滑线变阻器 $R$（其总阻值小于两平行电极间导电纸的电阻），
直流电压表 V（量程为 $6\UV$，内阻很大，其负接线柱与 B 极相连，正接线柱与探针 P 相连），
开关 K。
现要用图中仪器描绘两平行金属条 AB 间电场中的等势线。AB 间的电压要求取为 $6\UV$。
\begin{enumerate}
\item
在图中连线，画成实验电路原理图。
\item
下面是主要的实验操作步骤，将所缺的内容填写在横线上方。

a．接好实验电路。

b．\rule{4.0cm}{0.4pt}。

c．合上 K，并将探针 P 与 A 相接触。

d．\rule{4.0cm}{0.4pt}。

e．用探针压印的方法把 A、B 的位置标记在白纸上。画一线段连接 AB 两极，在连线上选取间距大致相等的 5 个点作为基准点，用探针把它们的位置印在白纸上。

f．将探针与某一基准点相接触，\rule{3.0cm}{0.4pt}，这一点是此基准的等势点。用探针把这一点的位置也压印在白纸上。用相同的方法找出此基准点的其它等势点。

g．重复步骤 f，找出其它 4 个基准点的等势点。取出白纸画出各条等势线。
\end{enumerate}
\onepicture[align=right, w=6.0cm]{figs/22b.png}
%% answer: 
\jdanswer{
\begin{enumerate}
\item
连接线路如图所示。
\item
b．把变阻器的滑动触头移到靠近 D 端处；d．调节 $R$，使电压表读数为 $6\UV$；f．记下电压表读数，在导电纸上移动探针，找出电压表读数与所记下的数值相同的另一点。
\end{enumerate}
}
%% memo:
\memoanswer{
（1）连接线路如图所示。

\onepicture[w=6.0cm]{figs/22b_an.jpg}

（2）b．把变阻器的滑动触头移到靠近 D 端处；d．调节 $R$，使电压表读数为 $6\UV$；f．记下电压表读数，在导电纸上移动探针，找出电压表读数与所记下的数值相同的另一点。
}

\item
%% number: 23
%% paperName: 2004年全国理综4第 23 题 16 分
%% typeId: 6
%% type: 计算题
%% chapter: 电路
%% point: 闭合电路欧姆定律
%% method: 无
%% score: 16
%% degree: 550
%% duplicateId: 0
%% body:
电动势为 $E=12\UV$ 的电源与一电压表和一电流表串联成闭合回路。如果将一电阻与电压表并联，则电压表的读数减小为原来的 $\frac{1}{3}$，电流表的读数增大为原来的 3 倍。求电压表原来的读数。
%% answer: 
\jdanswer{
$U=9\UV$
}
%% memo:
\memoanswer{
设电压表和电流表的原来读数分别为 $U$ 和 $I$，电源和电流表的内阻分别为 $r_{1}$ 和 $r_{2}$，由欧姆定律得
\[E=U+I(r_{1}+r_{2}) \tag{①}\]
\[E=\frac{1}{3}U+3I(r_{1}+r_{2}) \tag{②}\]
联立 ①② 并代入 $E$ 的数值得 $U=9\UV$。
}

\item
%% number: 24
%% paperName: 2004年全国理综4第 24 题 18 分
%% typeId: 6
%% type: 计算题
%% chapter: 磁场
%% point: 带电粒子在磁场中的运动
%% method: 无
%% score: 18
%% degree: 550
%% duplicateId: 0
%% body:
一匀强磁场，磁场方向垂直于 $xy$ 平面，在 $xy$ 平面上，磁场分布在以 O 为中心的一个圆形区域内。一个质量为 $m$、电荷量为 $q$ 的带电粒子，由原点 O 开始运动，初速为 $v$，方向沿 $x$ 正方向。后来，粒子经过 $y$ 轴上的 P 点，此时速度方向与 $y$ 轴的夹角为 $30^\circ$，P 到 O 的距离为 $L$，如图所示。不计重力的影响。求磁场的磁感强度 $B$ 的大小和 $xy$ 平面上磁场区域的半径 $R$。
\onepicture[align=right, w=5.5cm]{figs/24.png}
%% answer: 
\jdanswer{
$B=\frac{3mv}{qL}$，$R=\frac{\sqrt{3}}{3}L$
}
%% memo:
\memoanswer{
粒子在磁场中做匀速圆周运动，设其半径为 $r$，由洛伦兹力提供向心力得
\[qvB=m\frac{v^{2}}{r} \tag{①}\]
由题意知，粒子在磁场中的轨迹的圆心 C 必在 $y$ 轴上，且 P 点在磁场区之外，过 P 沿速度反向作延长线，它与 $x$ 轴相交于 Q 点，作圆弧过 O 点与 $x$ 轴相切，并且与 PQ 相切于 A，A 即粒子离开磁场区的点，这样可得圆弧轨迹的圆心 C，如图所示。
由图中几何关系得
\[L=3r \tag{②}\]
由 ①② 求得
\[B=\frac{3mv}{qL}\]
图中 OA 的长度即圆形磁场区的半径 $R$，由图中几何关系得
\[R=\frac{\sqrt{3}}{3}L\]
\onepicture[w=5.0cm]{figs/24_an.jpg}
}

\item
%% number: 25
%% paperName: 2004年全国理综4第 25 题 19 分
%% typeId: 6
%% type: 计算题
%% chapter: 运动和力
%% point: 动量和能量
%% method: 无
%% score: 19
%% degree: 450
%% duplicateId: 0
%% body:
如图，长木板 ab 的 b 端固定一档板，木板连同档板的质量为 $M=4.0\Ukg$，a、b 间距离 $s=2.0\Um$。木板位于光滑水平面上。在木板 a 端有一小物块，其质量 $m=1.0\Ukg$，小物块与木板间的动摩擦因数 $\mu=0.10$，它们都处于静止状态。现令小物块以初速 $v_{0}=4.0\Ums$ 沿木板向前滑动，直到和档板相撞。碰撞后，小物块恰好回到 a 端而不脱离木板。求碰撞过程中损失的机械能。
\onepicture[align=right, w=7.0cm]{\includesvg{figs/25}}
%% answer: 
\jdanswer{
$E_{1}=2.4\UJ$
}
%% memo:
\memoanswer{
设木板和物块最后共同的速度为 $v$，由动量守恒得
\[mv_{0}=(m+M)v \tag{①}\]
设全过程损失的机械能为 $E$，由能量守恒得
\[E=\frac{1}{2}mv_{0}^{2}-\frac{1}{2}(m+M)v^{2} \tag{②}\]
用 $s_{1}$ 表示物块从开始运动到碰撞前瞬间木板的位移，$W_{1}$ 表示在这段时间内摩擦力对木板所做的功，用 $W_{2}$ 表示同样时间内摩擦力对物块所做的功。
用 $s_{2}$ 表示从碰撞后瞬间到物块回到 a 端时木板的位移，$W_{3}$ 表示在这段时间内摩擦力对木板所做的功，用 $W_{4}$ 表示同样时间内摩擦力对物块所做的功。
用 $W$ 表示在全过程中摩擦力做的总功，则
\[W_{1}=\mu mgs_{1} \tag{③}\]
\[W_{2}=-\mu mg(s_{1}+s) \tag{④}\]
\[W_{3}=-\mu mgs_{2} \tag{⑤}\]
\[W_{4}=\mu mg(s_{2}-s) \tag{⑥}\]
\[W=W_{1}+W_{2}+W_{3}+W_{4} \tag{⑦}\]
用 $E_{1}$ 表示在碰撞过程中损失的机械能，则
\[E_{1}=E-W \tag{⑧}\]
由 ①至⑧ 式解得
\[E_{1}=\frac{1}{2}\cdot\frac{mM}{m+M}v_{0}^{2}-2\mu mgs\]
代入数据得 $E_{1}=2.4\UJ$。
}

\end{enumerate}

\end{document}
